\documentclass{article}
\usepackage[utf8]{inputenc}
\usepackage[brazil]{babel}

\title{aula 3 - respostas}
\author{marina}
\date{\today}

\begin{document}

\maketitle

\section{explicação masiero}
\begin{itemize}
    \item grande concentração em baixas frequências -> som ambiente
    \item linhas verticais -> energia em todas as frequências -> dias com muito vento -> vento bate no microfone -> faz ruído
    \item linhas horizontais -> ruído elétrico -> frequência da rede 60hz 
    \item fontes do sistema elétrico vibram em ressonância com a rede
    \item ar absorve mais agudo do que grave
\end{itemize}

\section{exercício 04}
usos de engenharia de som além da música
\begin{itemize}
    \item cancelar harmônico mais proeminente - controle ativo de ruído - gerar sinal com fase oposta pra cancelar
    \item monitoramento de máquinas -> manutenção preventiva
    \item aparelho auditivo -> aasi
\end{itemize}

\section{exercício 05}
\begin{itemerize}
    \item a partir de 10ms já diminui bastante,a partir de 7, já fica só um clique
    \item o mais grave parece mais baixo
    \item 
\end{itemerize}
\end{document}
